\silentchapter{Til lesaren}{til lesaren}
\begin{widebody}
\begin{foreordbody}

Boka har sju delar.

\textit{Ordlista} definerer dei tekniske termane i alfabetisk orden. Ho er til for å slå opp i undervegs, ikkje for å lesast fyrst.

\textit{Traktaten} er aksiomatisk. Ho byggjer eit omgrepsapparat frå grunnen av, gjennom nummererte proposisjonar i åtte seksjonar. Siste seksjon er den generative modellen, der grammatikk, landskap og lyskjegle blir sameinte i éi likning. Proposisjonane er merkte etter type: definisjon (\textit{d}), teorem (\textit{t}), postulat (\textit{a}), observasjon (\textit{o}), illustrasjon (\textit{i}).

\textit{Fundamentet} viser at rammeverket ikkje er isolert. Same formalismen er empirisk demonstrert over fjorten substratnivå, frå genregulatoriske nettverk til sivilisasjonar. Forskingsprogrammet som følgjer står skissert i fire punkt.

\textit{Etterord I} er materiellt og visceralt. Han skildrar kva form sviktar som, når omhugen er for trong, i ei lang rad eksempel frå eingongskoppen til overvakingsarkitekturen.

\textit{Etterord II} er programmatisk. Han stiller kravet til designfaget direkte.

\textit{Referansane} er nummererte. Superscript-tala i proposisjonane og etterorda peikar hit.

\textit{Notata} kommenterer enkeltproposisjonar, plasserer dei i tradisjonen, og listar falsifiseringsvilkåra.

Akademikaren vil finne det meste i traktaten, fundamentet og notata. Studenten vil finne seg sjølv i etterorda. Designaren i praksis bør lese siste sida fyrst, og byrje på nytt.

\end{foreordbody}
\end{widebody}
